\subsection{Local normalization at the chain and a conductor-degree obstruction}

\begin{proposition}[P-F75: a conductor generator must detect the exponential stabilizer]
\label{prop:localconductor}
Use the rings $R\twoheadrightarrow A\subset B$, the grading,
and the group $G'$ of P-F71, with $n\ge4$.
Let $T_n\in S=\operatorname{Sym}^3\mathbb C^n$ be the
balanced chain tensor of P-F60, and let $q_n\in\operatorname{Spec}B$
and $p_n\in\operatorname{Spec}A\subset\operatorname{Spec}R$
be its images.  Write $\mathfrak m_n\subset A$ for the
maximal ideal of $p_n$, and set
\[
 I_n=\{a\in A:aB\subset A\},\qquad M=2^{n-1}.
                                                               \tag{P-F75a}
\]
Then the finite normalization in P-F71 is an isomorphism
near $p_n$.  In particular
\[
 A_{\mathfrak m_n}=B_{\mathfrak m_n},\qquad
 (I_n)_{\mathfrak m_n}=A_{\mathfrak m_n}.
                                                               \tag{P-F75b}
\]
Here $B_{\mathfrak m_n}$ means localization as a finite
$A$-algebra; it is local because there is exactly one point
over $p_n$.

There is a homogeneous conductor element $h\in I_n$ with
$h(q_n)\ne0$.  Every homogeneous element of $B$ nonzero
at $q_n$ has degree divisible by $M$, in the grading
$B_s=\mathbb C[S]_{ns}^{SL_n}$.  Consequently, on the
explicit nonsurjective family
\[
 n=\binom{3r-1}{2},\qquad r\ge2,
                                                               \tag{P-F75c}
\]
every homogeneous conductor element nonzero at $q_n$
has degree at least $2^{n-1}$, and every homogeneous
generating set of the conductor ideal $I_n$ contains
an element of degree at least $2^{n-1}$.
This is a lower bound on conductor elements which do
not vanish at the chain, and on at least one conductor
generator.  It is not a lower bound on every nonzero
element of the conductor.
\end{proposition}

\begin{proof}
\emph{The effective stabilizers and tangent spaces.}
Put $V=\mathbb C^n$, $W=V^{\otimes3}$, and
$H=SL(V)^3$.  P-F60 proves that $T_n$ has a closed
$H$-orbit, that its stabilizer in $H$ is exactly the
ineffective central kernel, and that its centroid is
\[
 \{L\in\operatorname{End}(V):(L\otimes I\otimes I)T_n
                  =(I\otimes L\otimes I)T_n\}=\mathbb C I.
                                                               \tag{P-F75d}
\]
We also need the effective stabilizer for the exact
diagonal group $G'=\{D:(\det D)^3=1\}$, rather than
for an unspecified general linear action.
If $D\in G'$ fixes $T_n$, then $(D,D,D)$ belongs to
the full product stabilizer (P-F60c), and its product
determinant character is $(\det D)^3=1$.
The computation (P-F60d) forces its dyadic parameter
to equal one.  Hence $D=zI$, and fixing the nonzero
cubic tensor means $z^3=1$.  Conversely every such
scalar belongs to $G'$ and acts trivially on all of $S$.
There are no further elements in the ineffective kernel:
if $D$ fixes every $v^{\otimes3}$, then it preserves
every line in $V$, hence is scalar, and that scalar
has cube one.  Thus
\[
 \operatorname{Stab}_{G'}(T_n)=\boldsymbol\mu_3 I
               =\ker(G'\curvearrowright S).
                                                               \tag{P-F75e}
\]
The diagonal orbit is closed, either by the equal
marginal-moment computation in P-F71 or directly by
the balance of P-F60.

The free effective slice at these closed orbits gives
smooth quotient points and the natural tangent
identifications
\[
 T_{q_n}\operatorname{Spec}B
       =S/(\mathfrak{sl}(V)\cdot T_n),\qquad
 T_{p_n}\operatorname{Spec}R
       =W/(\mathfrak{sl}(V)^3\cdot T_n).
                                                               \tag{P-F75f}
\]
The diagonal action is meant in the first denominator.
The primary slice statement used is Luna,
\emph{Slices \'etales}, Section~III.1, theorem on
printed p.~97 and Remark~$1^\circ$ on pp.~97--98:
\url{https://www.numdam.org/item/10.24033/msmf.110.pdf}.
One can apply its connected-group version first to
$SL(V)$ and then take the finite residual quotient:
the stabilizer in $SL(V)$ is the ineffective scalar
group $\boldsymbol\mu_{\gcd(n,3)}$, while
$G'/(SL(V)\boldsymbol\mu_3)$ acts freely at the
resulting quotient point by (P-F75e) and uniqueness
of closed orbits in a quotient fiber.
A finite-group quotient with trivial point stabilizer
is \'etale at that point in characteristic zero, so
it preserves the tangent space and smoothness in
(P-F75f).  Explicitly, completion over the finite orbit
decomposes the finite algebra as the product of the
completed local rings at the orbit points.  Invariants
commute with this completion, since averaging is an
idempotent linear operator in characteristic zero.
When the stabilizer is trivial, invariants of that
permuted product identify with any one factor.
Thus the completed local ring of the quotient equals
the completed local ring at the chosen point, which
also directly proves the tangent and smoothness assertions.
This handles $3\mid n$ without
silently replacing $G'$ by its identity component.

\emph{The differential of restriction is injective.}
The map from the first tangent space in (P-F75f) to
the second is induced by the actual inclusion $S\subset W$.
We prove
\[
 S\cap(\mathfrak{sl}(V)^3\cdot T_n)
                  =\mathfrak{sl}(V)\cdot T_n.
                                                               \tag{P-F75g}
\]
Suppose $X,Y,Z\in\mathfrak{sl}(V)$ and
\[
 u=(X\otimes I\otimes I+I\otimes Y\otimes I
                         +I\otimes I\otimes Z)T_n\in S.
\]
Both $u$ and $T_n$ are symmetric in the first two
positions.  Subtracting the identity obtained by
interchanging those positions gives
\[
 ((X-Y)\otimes I\otimes I)T_n
                    =(I\otimes(X-Y)\otimes I)T_n.
\]
By (P-F75d), $X-Y$ is scalar; since it is traceless,
it vanishes.  Interchanging the first and third
positions gives $X-Z=0$ by the same centroid identity,
using symmetry to permute positions.  Therefore
$X=Y=Z$, proving (P-F75g).  The reverse inclusion is
immediate.  It follows that
\[
 T_{q_n}\operatorname{Spec}B
                \longrightarrow T_{p_n}\operatorname{Spec}R
                     \quad\hbox{is injective}.
                                                               \tag{P-F75h}
\]
Since this morphism factors through the closed image
$\operatorname{Spec}A$, its differential to
$T_{p_n}\operatorname{Spec}A$ is injective as well.

\emph{The finite fiber and the local isomorphism.}
P-F71 already proves that $A\subset B$ is finite and
that there is exactly one closed point of
$\operatorname{Spec}B$ over $p_n$.
Let $A'=A_{\mathfrak m_n}$ and $B'=B_{\mathfrak m_n}$.
The algebra $B'$ is finite and local over $A'$, and
its fiber
\[
                         C=B'/\mathfrak m_nB'
                                                               \tag{P-F75i}
\]
is an Artinian local $\mathbb C$-algebra with residue
field $\mathbb C$.
Its Zariski tangent space is the kernel of the
differential $T_{q_n}\operatorname{Spec}B\to
T_{p_n}\operatorname{Spec}A$: this follows directly
from the definition of tangent vectors as maps from
$\mathbb C[\epsilon]/(\epsilon^2)$ with fixed base point.
It is zero by (P-F75h).  If $\mathfrak r$ is the
maximal ideal of $C$, then
$\mathfrak r/\mathfrak r^2=0$; Nakayama's lemma gives
$\mathfrak r=0$.  Thus $C=\mathbb C$, generated by
the image of $1$.

Applying Nakayama again, now to the finite $A'$-module
$B'/A'\cdot1$, shows that it is zero.
As $A'\subset B'$ is an inclusion, $A'=B'$.
Equivalently the finite $A$-module $B/A$ vanishes
after localization at $\mathfrak m_n$.
Since
\[
                  I_n=\operatorname{Ann}_A(B/A),
                                                               \tag{P-F75j}
\]
some element of $I_n$ is outside $\mathfrak m_n$:
choose finitely many generators of $B/A$ and multiply
their individual annihilators outside $\mathfrak m_n$.
Inverting this product also gives an open neighborhood
on which $A=B$, proving the local-isomorphism assertion.
The conductor is graded, as proved in P-F71.
If a conductor element has nonzero value at $q_n$,
at least one of its homogeneous components has nonzero
value there; that component is still in $I_n$.
This proves the existence of the required homogeneous $h$.

\emph{The exact grading and the lower bound.}
Restriction respects ordinary degree and the common
stretch grading $ns$.  Hence the finite bijection
$\operatorname{Spec}B\to\operatorname{Spec}A$ is
equivariant for the grading $\mathbb C^*$ action.
Its injectivity on closed points implies that the
grading stabilizers of $q_n$ and $p_n$ coincide.
The latter is exactly $\boldsymbol\mu_M$ by P-F60;
the closed immersion into $\operatorname{Spec}R$
does not change that stabilizer.  Consequently
\[
 \operatorname{Stab}_{\mathbb C^*}(q_n)=\boldsymbol\mu_M,
 \qquad h\in B_c, h(q_n)\ne0\ \Longrightarrow\ M\mid c.
                                                               \tag{P-F75k}
\]
Indeed evaluation at a primitive $M$th root $\zeta$
gives $h(q_n)=h(\zeta q_n)=\zeta^c h(q_n)$.
This uses ordinary degree $nc$, not the full cubic
canonical unit $n/\gcd(n,3)$; the finite bijection
therefore verifies the exact value of $M$ after the
Veronese regrading, rather than assuming it is unchanged.

For (P-F75c), P-F71m--q proves $A\ne B$.
Thus the conductor contains no nonzero degree-zero
element, since such an element would be an invertible
scalar and imply $B=A$.
Every homogeneous $h\in I_n$ nonzero at $q_n$ therefore
has positive degree divisible by $M$, giving $\deg h\ge M$.
Finally let $h_1,\ldots,h_v$ be any homogeneous generating
set of $I_n$.  If all $h_j$ vanished at $q_n$, every
element of $I_n$ would vanish there, contradicting
(P-F75b).  At least one $h_j$ is nonzero there and
has degree at least $M$.  This proves the stated
generator bound.
\end{proof}

The result quantifies one obstruction to clearing the
normalization denominators while retaining the chain
point.  It does not replace the unspecified conductor
shift in (P-F71d) by an exponential lower bound: a
different homogeneous conductor element may vanish
at the chain.  It also gives no global tensor Hilbert
pole or strict Kronecker-index conclusion.
